% GENERATED FILE. Edit canonical lessons, then run npm run generate:books.
% canonical-source-hash: fnv1a64:920689d8176f396b
% canonical-lessons: ZH-C255-dagong, ZH-W255-gong-work, ZH-C255-gongzi, ZH-C255-tongshi, ZH-C255-jiaban, ZH-C255-laoban

\chapter{A Part-Time Job}
\label{ch:zh-a-part-time-job}
\hlchaptermodality{255}

\begin{chapteropening}
\textbf{By the end of this chapter:} \emph{I can talk about a part-time job, pay, colleagues, working late and the boss.}

Complete the last lesson of chapter 255: I can talk about a part-time job, pay, colleagues, working late and the boss.
\end{chapteropening}

\section[\texorpdfstring{\zh{打工} (dǎgōng)}{打工 (dǎgōng)}]{\zh{打工} (dǎgōng) — to work part-time; to have a part-time job}

\label{lesson:ZH-C255-dagong}



\begin{quote}
Before the new one: say the Mandarin for a wedding, then the Mandarin for daily life.
\end{quote}

\subsection*{You'll want to know: \zh{打工}}
\textbf{\zh{打工}} — \emph{dǎgōng} — \textquotedblleft{}to work part-time; to have a part-time job\textquotedblright{}.

\begin{quote}
\textbf{\zh{他晚上打工}} \emph{tā wǎnshang dǎgōng} — He works a part-time job in the evenings.
\end{quote}

The new character, in pieces: \zh{工} — three strokes: a short bar, a vertical down from its middle, and a long bar across the foot. The next lesson puts it in your hand.

\begin{cousinweb}
\textbf{\zh{打工}} (\emph{dǎgōng}) is built from \zh{打} 'to strike' and \zh{工} 'work', both inherited from Middle Chinese, \emph{taengX} and \emph{kuwng} (Guangyun, in Baxter's transcription), which glosses the first 'to strike' and the second 'skill'. Middle Chinese \emph{taengX} would regularly give dǎng; dǎ is irregular. How the pair came to mean casual work is not settled.
\end{cousinweb}

\begin{grammarlens}[title={What you've built}]
One more word for saying what I do.
\end{grammarlens}

\subsection*{Your turn}
\begin{itemize}
\raggedright
  \item \emph{Say it:} \emph{dǎgōng}
  \item \emph{Say it:} \emph{dǎgōng}, once more
  \item \emph{Say it:} \emph{tā wǎnshang dǎgōng} — He works a part-time job in the evenings.
  \item \emph{Recall:} say the Mandarin for a wedding, then the Mandarin for daily life, then say \emph{dǎgōng} again
\end{itemize}

\begin{tcolorbox}[breakable,colback=teal!4,colframe=teal!35!black,title={Before you move on}]
What does \zh{打工} mean? (\textquotedblleft{}To work part-time; to have a part-time job\textquotedblright{}.) Say it once more.
\end{tcolorbox}
\section[\texorpdfstring{\zh{工} (gōng)}{工 (gōng)}]{\zh{工} — two bars and a post}

\label{lesson:ZH-W255-gong-work}



\begin{quote}
The word first, then the character inside it.

\begin{itemize}
\raggedright
  \item \emph{Recall:} say \emph{dǎgōng} — \textbf{R1}, one lesson back, before it decays
\end{itemize}

\begin{quote}
\zh{在} \zh{打}
\end{quote}

\emph{Read it:} each one aloud, with nothing said first — \emph{zài}, \emph{dǎ}
\end{quote}

\subsection*{Script — two bars and a post}
\hlblockfigure{\detokenize{figures/ZH-W255-gong-work-filmstrip.pdf}}{How \zh{工} is written, stroke by stroke}

\begin{quote}
\zh{工}
\end{quote}

\textbf{One piece, three strokes, top to bottom.}

\textbf{A short bar} across the top.

\textbf{A vertical} down from the middle of the bar. It starts at the bar and does not cross it.

\textbf{A long bar} across the foot, longer than the top one, closing the vertical.

\subsection*{Your turn}
\begin{itemize}
\raggedright
  \item \emph{Write it:} \zh{工}, all three, with the model in front of you
  \item \emph{Cover:} write it again from memory, and count the pieces, not the strokes
\end{itemize}

\begin{tcolorbox}[breakable,colback=teal!4,colframe=teal!35!black,title={Before you move on}]
How many strokes in \zh{工}? (\textbf{Three}.) Which bar is longer? (\textbf{The one at the foot}.) Does the vertical cross the top bar? (\textbf{No, it starts there}.)

Source: \href{https://raw.githubusercontent.com/chanind/hanzi-writer-data/68d10a4b21150cae5e1ebbd223eed289cf32d90c/data/\%E5\%B7\%A5.json}{Hanzi Writer Data \zh{工}, pinned stroke paths and medians}.
\end{tcolorbox}
\section[gōngzī]{gōngzī — pay; wages}

\label{lesson:ZH-C255-gongzi}



\begin{quote}
Before the new one: say the Mandarin for daily life, then the Mandarin for to work part-time.
\end{quote}

\subsection*{You'll want to know: gōngzī}
\emph{gōngzī} — \textquotedblleft{}pay; wages\textquotedblright{}.

\begin{quote}
\emph{zhège gōngzuò de gōngzī hěn gāo} — This job pays well.
\end{quote}

Its characters come later; for now it is a word for the ear.

\begin{cousinweb}
\emph{gōngzī} is built from \zh{工} 'work' and \emph{zī} 'goods', both inherited from Middle Chinese, \emph{kuwng} and \emph{tsij} (Guangyun, in Baxter's transcription), which glosses the second 'goods' and 'to help'.
\end{cousinweb}

\begin{grammarlens}[title={What you've built}]
One more word for saying what I do.
\end{grammarlens}

\subsection*{Your turn}
\begin{itemize}
\raggedright
  \item \emph{Say it:} \emph{gōngzī}
  \item \emph{Say it:} \emph{gōngzī}, once more
  \item \emph{Say it:} \emph{zhège gōngzuò de gōngzī hěn gāo} — This job pays well.
  \item \emph{Recall:} say the Mandarin for daily life, then the Mandarin for to work part-time, then say \emph{gōngzī} again
\end{itemize}

\begin{tcolorbox}[breakable,colback=teal!4,colframe=teal!35!black,title={Before you move on}]
What does gōngzī mean? (\textquotedblleft{}Pay; wages\textquotedblright{}.) Say it once more.
\end{tcolorbox}
\section[\texorpdfstring{\zh{同事} (tóngshì)}{同事 (tóngshì)}]{\zh{同事} (tóngshì) — colleague}

\label{lesson:ZH-C255-tongshi}



\begin{quote}
Before the new one: say the Mandarin for to work part-time, then the Mandarin for wages.
\end{quote}

\subsection*{You'll want to know: \zh{同事}}
\textbf{\zh{同事}} — \emph{tóngshì} — \textquotedblleft{}colleague\textquotedblright{}.

\begin{quote}
\textbf{\zh{他是我的同事我们一起工作}} \emph{tā shì wǒ de tóngshì wǒmen yìqǐ gōngzuò} — He's a colleague of mine; we work together.
\end{quote}

\begin{cousinweb}
\textbf{\zh{同事}} (\emph{tóngshì}) is built from \zh{同} 'the same' and \zh{事} 'affairs', both inherited from Middle Chinese, \emph{duwng} and \emph{dzriH} (Guangyun, in Baxter's transcription), which glosses the first 'together': people who share the work.
\end{cousinweb}

\begin{grammarlens}[title={What you've built}]
One more word for saying what I do.
\end{grammarlens}

\subsection*{Your turn}
\begin{itemize}
\raggedright
  \item \emph{Say it:} \emph{tóngshì}
  \item \emph{Say it:} \emph{tóngshì}, once more
  \item \emph{Say it:} \emph{tā shì wǒ de tóngshì wǒmen yìqǐ gōngzuò} — He's a colleague of mine; we work together.
  \item \emph{Recall:} say the Mandarin for to work part-time, then the Mandarin for wages, then say \emph{tóngshì} again
\end{itemize}

\begin{tcolorbox}[breakable,colback=teal!4,colframe=teal!35!black,title={Before you move on}]
What does \zh{同事} mean? (\textquotedblleft{}Colleague\textquotedblright{}.) Say it once more.
\end{tcolorbox}
\section[jiābān]{jiābān — to work overtime}

\label{lesson:ZH-C255-jiaban}



\begin{quote}
Before the new one: say the Mandarin for wages, then the Mandarin for a colleague.
\end{quote}

\subsection*{You'll want to know: jiābān}
\emph{jiābān} — \textquotedblleft{}to work overtime\textquotedblright{}.

\begin{quote}
\emph{jīntiān wǒ yào jiābān} — I have to work late today.
\end{quote}

Its characters come later; for now it is a word for the ear.

\begin{cousinweb}
\emph{jiābān} is built from \emph{jiā} 'to add' and \emph{bān}, both inherited from Middle Chinese, \emph{kae} and \emph{paen} (Guangyun, in Baxter's transcription), which glosses the first 'to increase' and quotes an older dictionary for the second, 'to divide jade'. The sense 'a shift at work' is not given there, so how it came is not settled.
\end{cousinweb}

\begin{grammarlens}[title={What you've built}]
One more word for saying what I do.
\end{grammarlens}

\subsection*{Your turn}
\begin{itemize}
\raggedright
  \item \emph{Say it:} \emph{jiābān}
  \item \emph{Say it:} \emph{jiābān}, once more
  \item \emph{Say it:} \emph{jīntiān wǒ yào jiābān} — I have to work late today.
  \item \emph{Recall:} say the Mandarin for wages, then the Mandarin for a colleague, then say \emph{jiābān} again
\end{itemize}

\begin{tcolorbox}[breakable,colback=teal!4,colframe=teal!35!black,title={Before you move on}]
What does jiābān mean? (\textquotedblleft{}To work overtime\textquotedblright{}.) Say it once more.
\end{tcolorbox}
\section[lǎobǎn]{lǎobǎn — boss; owner (of a business)}

\label{lesson:ZH-C255-laoban}



\begin{quote}
Before the new one: say the Mandarin for a colleague, then the Mandarin for to work overtime.
\end{quote}

\subsection*{You'll want to know: lǎobǎn}
\emph{lǎobǎn} — \textquotedblleft{}boss; owner (of a business)\textquotedblright{}.

\begin{quote}
\emph{wǒ de lǎobǎn rén hěn hǎo} — My boss is a nice person.
\end{quote}

Its characters come later; for now it is a word for the ear.

Two third tones meet in this word: lǎo is said with a rising tone before bǎn, and the spelling keeps both marks.

\begin{cousinweb}
\emph{lǎobǎn} is built from \emph{lǎo} 'old' and \emph{bǎn} 'a board', both inherited from Middle Chinese, \emph{lawX} and \emph{paenX} (Guangyun, in Baxter's transcription). How the pair came to mean the boss is not settled.
\end{cousinweb}

\begin{grammarlens}[title={What you've built}]
One more word for saying what I do.
\end{grammarlens}

\subsection*{Your turn}
\begin{itemize}
\raggedright
  \item \emph{Say it:} \emph{lǎobǎn}
  \item \emph{Say it:} \emph{lǎobǎn}, once more
  \item \emph{Say it:} \emph{wǒ de lǎobǎn rén hěn hǎo} — My boss is a nice person.
  \item \emph{Recall:} say the Mandarin for a colleague, then the Mandarin for to work overtime, then say \emph{lǎobǎn} again
\end{itemize}

\begin{tcolorbox}[breakable,colback=teal!4,colframe=teal!35!black,title={Before you move on}]
What does lǎobǎn mean? (\textquotedblleft{}Boss; owner (of a business)\textquotedblright{}.) Say it once more.
\end{tcolorbox}
